\documentclass[10pt,a4paper]{amsart}
\usepackage[margin=19mm]{geometry}
\usepackage{amsmath,amssymb,mathtools}
\usepackage[scr=rsfs,cal=euler]{mathalfa}
\usepackage{xfrac}
\usepackage[unicode,hidelinks,bookmarks=false]{hyperref}
\newcommand{\ie}{i.e.\ }
\newcommand{\cf}{cf.\ }
\newcommand{\resp}{respectively\ }
\newcommand{\Subsection}{\subsection}
\providecommand{\nobreakdash}{\nobreak\hbox{-}}
\setlength{\emergencystretch}{2em}
\multlinegap0pt
\usepackage{xfrac}
\usepackage{amssymb}
\usepackage{mathtools}
\usepackage{tcolorbox}
\usepackage[normalem]{ulem}
\newtheorem{prop}{Proposition}
\newcommand{\R}{\mathbb R}
\newcommand{\be}{\begin{equation}}
\newcommand{\ee}{\end{equation}}
\newcommand{\na}{\nabla}
\title{On a Multispecies Electrodiffusion Model in Higher-Dimensional Porous Media}
\author{Elie Abdo, Christofer Ghazale}
\date{Source: arXiv:2609.14659v1 (CC BY 4.0)}
\begin{document}
\maketitle

\noindent\emph{Source and licence.} On a Multispecies Electrodiffusion Model in Higher-Dimensional Porous Media. Version arXiv:2609.14659v1 (CC BY 4.0), distributed by arXiv under the Creative Commons Attribution 4.0 International licence, http://creativecommons.org/licenses/by/4.0/, as recorded in the arXiv licence metadata for this exact version and shown on the abstract page https://arxiv.org/abs/2609.14659v1. Full licence notice: \texttt{licenses/arxiv-2609.14659-CC-BY-4.0.txt}.\par\smallskip

\noindent\textit{Editorial scope.} Counted unit: the article's prop prop2.1 (source0.tex lines 271--276). The article's complete original proof of it is delivered with it (source0.tex lines 278--290, one \texttt{proof} environment of 13 lines). No mathematics is written, repaired or re-derived: the statement and the proof are the article's own lines, in the article's own notation, and no retained line is substituted or re-flowed.  No further source-local context is needed: every cross-reference of every retained line resolves inside the two delivered spans.  Nothing was omitted from any retained span, so the package carries no omission record.  No retained line contains a citation, so the wrapper carries no bibliography and no bibliography entry.  No retained line contains a figure environment, an image-inclusion command or drawing code, so no image is delivered and the package declares no asset dependency.  Editorial formatting only, in addition: an AMS \texttt{amsart} preamble that restates the article's own declarations and macros used by the retained lines, namely the environments \texttt{prop} on separate counters, together with the standard packages those lines need.  The retained environments print in this document as Proposition 1 in order of appearance, whereas the article prints the same items with the numbering of its own full document.  The complete native source of the article as distributed in the arXiv e-print for this exact version is bundled unchanged as \texttt{sources/210343/source0.tex}.
\begin{prop}\label{prop2.1}
    Let $d \ge 3$ be an integer, and let $p > d$ be a real number. Let $g$ be in $L^1(\R^d) \cap L^p(\R^d)$ such that $\int_{\R^d} g = 0$. Let $\Psi$ be a solution to the Poisson equation $- \Delta \Psi = g$  with decay at infinity. Then the following elliptic regularity estimate
    \be
\|\na \Psi\|_{L^{\infty}} \le C\left(\|g\|_{L^1} + \|g\|_{L^p} \right)
    \ee holds.
\end{prop}

\begin{proof}
    We have
    \be 
\begin{aligned}
\|\nabla\Psi\|_{L^\infty} 
&\le C(\|\nabla\Psi\|_{L^p} +\|g\|_{L^p}) 
\le C(\|(-\Delta)^{-1/2}g\|_{L^p} + \|g\|_{L^p}) 
\le C(\|g\|_{L^\frac{pd}{d+p}} + \|g\|_{L^p})
\\&\le C(\|g\|_{L^1}^\theta \|g\|_{L^p}^{1-\theta} + \|g\|_{L^p}) 
\le C\|g\|_{L^1} + C\|g\|_{L^p}
\end{aligned}
\ee where $\theta \in (0,1)$ is the parameter resulting from interpolating $L^{\frac{pd}{d+p}}$ between $L^1$ and $L^p$. Here, we made use of the Morrey, Hardy-Littlewood-Sobolev, and Young inequalities.
\end{proof}

\end{document}
